\section{Experiments}
\label{sec:experiments}

In this section, we present a rigorous empirical evaluation of our proposed routers against standard and state-of-the-art model merging baselines on fully converged expert networks. Our experiments are designed to systematically test our methodological hypotheses and deconstruct the performance of quantum wavefunction superposition.

\subsection{Experimental Setup}
\textbf{Tasks and Datasets:} We evaluate all methods on a diverse, high-conflict multi-task suite comprising four image classification tasks: MNIST, FashionMNIST, CIFAR-10, and the Street View House Numbers (SVHN) dataset. This selection spans grayscale digits, grayscale fashion items, low-resolution natural objects, and complex colored digits, presenting severe gradient conflicts in weight space.

\textbf{Backbone and Experts:} We use the compact Vision Transformer (\texttt{vit\_tiny\_patch16\_224}, 5.7M parameters) as the base model. Training is conducted with a hidden dimension $D = 192$ and 12 transformer blocks. This capacity-constrained backbone is highly susceptible to parameter interference, providing an ideal testbed for weight routing. To resolve previous experimental pipeline flaws, we fine-tune four task-specific experts to high convergence using AdamW on GPU with a dual learning rate schedule ($2\times 10^{-5}$ for the backbone, $10^{-3}$ for heads) for exactly 15 epochs. Each specialized expert achieves high converged accuracy (100.00\% on MNIST, 92.80\% on FashionMNIST, 96.40\% on CIFAR-10, and 96.80\% on SVHN), establishing a strong upper bound.

\textbf{Calibration Protocol:} We use a balanced calibration set of $64$ total samples ($16$ samples per task). All trainable methods (OFS-Tune, Linear Router, QWS-Merge, BL-Router, and GLS-Router) are optimized on this calibration set for exactly $100$ steps of Adam with a learning rate of $1\times 10^{-2}$.

\textbf{Evaluation Protocol:} We perform two types of evaluations:
\begin{enumerate}
    \item \textbf{Homogeneous Evaluation:} Measures the joint multi-task capability on $1000$ independent, randomly sampled test set images per task.
    \item \textbf{Heterogeneous Evaluation:} Measures dynamic adaptation under streaming noise on a fully shuffled, interleaved stream of $2000$ test samples ($500$ per task) across batch sizes $B \in \{1, 16, 256\}$.
\end{enumerate}

\subsection{Homogeneous Multi-Task Performance}
Table~\ref{tab:homogeneous_results} presents the homogeneous multi-task joint accuracy for all evaluated models on converged task experts.

\begin{table*}[t]
\centering
\caption{Homogeneous multi-task joint accuracy (\%) on independent test sets (Mean $\pm$ Std over 3 random calibration-sampling seeds). Bold indicates the best-performing dynamic merging or routing method.}
\label{tab:homogeneous_results}
\vskip 0.15in
\begin{tabular}{lccccc}
\toprule
Method & MNIST & FashionMNIST & CIFAR-10 & SVHN & Joint Mean \\
\midrule
Individual Experts (Ceiling) & $100.00 \pm 0.00$\% & $92.80 \pm 0.00$\% & $96.40 \pm 0.00$\% & $96.80 \pm 0.00$\% & $96.50 \pm 0.00$\% \\
\midrule
Uniform Merge (Task Arithmetic) & $93.60 \pm 0.00$\% & $75.60 \pm 0.00$\% & $93.60 \pm 0.00$\% & $77.60 \pm 0.00$\% & $85.10 \pm 0.00$\% \\
AdaMerging (Test-Time Adaptation) & $96.40 \pm 0.00$\% & $81.60 \pm 0.00$\% & $92.40 \pm 0.00$\% & $86.80 \pm 0.00$\% & \textbf{89.30} $\pm$ \textbf{0.00}\% \\
OFS-Tune (Supervised Static) & $93.33 \pm 2.93$\% & $80.53 \pm 0.82$\% & $87.60 \pm 2.59$\% & $87.73 \pm 1.64$\% & $87.30 \pm 1.57$\% \\
\midrule
Linear Router (Classical Baseline) & $91.20 \pm 4.49$\% & $66.53 \pm 2.64$\% & $80.80 \pm 7.72$\% & $74.00 \pm 16.14$\% & $78.13 \pm 2.54$\% \\
Linear Router (\textbf{Reg - Ours}) & $89.60 \pm 11.92$\% & $65.87 \pm 6.71$\% & $84.00 \pm 1.50$\% & \textbf{91.73} $\pm$ \textbf{3.71}\% & $82.80 \pm 4.85$\% \\
QWS-Merge (SOTA Waveform) & $92.53 \pm 1.15$\% & $74.93 \pm 1.61$\% & $88.67 \pm 2.00$\% & $79.73 \pm 0.75$\% & $83.97 \pm 0.53$\% \\
\midrule
\textbf{BL-Router (Ours - Bounded Linear)} & $80.80 \pm 16.40$\% & $64.93 \pm 0.94$\% & $87.20 \pm 2.26$\% & $31.73 \pm 16.03$\% & $66.17 \pm 0.24$\% \\
\textbf{BL-Router (Ours - Reg)} & $73.47 \pm 13.97$\% & $64.93 \pm 1.00$\% & $90.80 \pm 0.33$\% & $43.20 \pm 8.02$\% & $68.10 \pm 3.07$\% \\
\textbf{GLS-Router (Ours - Layer Scaling)} & $87.60 \pm 8.45$\% & $64.80 \pm 3.53$\% & $82.27 \pm 1.91$\% & $74.67 \pm 24.30$\% & $77.33 \pm 6.71$\% \\
\textbf{GLS-Router (Ours - Reg)} & $87.73 \pm 9.64$\% & $64.67 \pm 3.03$\% & $83.73 \pm 4.69$\% & $78.53 \pm 20.54$\% & $78.67 \pm 6.81$\% \\
\textbf{BSigmoid-Router (Ours - Sigmoidal)} & $93.87 \pm 0.38$\% & $73.60 \pm 2.83$\% & $92.13 \pm 2.07$\% & $75.33 \pm 3.21$\% & $83.73 \pm 1.93$\% \\
\textbf{BSigmoid-Router (Ours - Reg)} & $93.87 \pm 0.38$\% & $74.00 \pm 2.26$\% & $91.87 \pm 2.45$\% & $74.80 \pm 3.96$\% & $83.63 \pm 2.07$\% \\
\bottomrule
\end{tabular}
\end{table*}

\begin{figure}[t]
\centering
\includegraphics[width=0.48\textwidth]{comparison_plot.png}
\caption{Homogeneous joint mean accuracy comparison across all models. Once expert models are fully converged, the performance gaps and structural profiles of all methods are accurately evaluated.}
\label{fig:comparison_plot}
\end{figure}

The empirical findings on fully converged expert models in Table~\ref{tab:homogeneous_results} yield several critical methodological insights:
\begin{enumerate}
    \item \textbf{Paradigm Clarification (TTA vs.\ Offline Calibration):} Under a strictly fair evaluation pipeline, QWS-Merge achieves a joint multi-task mean accuracy of \textbf{83.97 $\pm$ 0.53\%}, while the unsupervised test-time adaptation (TTA) baseline AdaMerging achieves \textbf{89.30 $\pm$ 0.00\%}. However, we highlight that comparing them purely on accuracy conflates two fundamentally different operational paradigms. AdaMerging dynamically optimizes merging coefficients on-the-fly during inference via active gradient updates to minimize prediction entropy, which introduces significant latency, backward-pass overhead, and susceptibility to stream noise. In contrast, offline-calibrated dynamic routing (like QWS-Merge) is trained once on a tiny offline calibration set and operates as a pure forward pass during inference with zero test-time active optimization or latency overhead, making it operationally far more practical for real-time deployment.
    \item \textbf{The Impact of Routing Head Regularization \& The Over-Scaling Confounder:} The standard unregularized Linear Router suffers from representational collapse on SVHN ($74.00 \pm 16.14$\%) and high instability due to overfitting on the tiny calibration set. Strikingly, when we apply standard L2 regularization during calibration, our \textbf{Linear Router (Reg)} completely resolves the SVHN collapse, boosting SVHN accuracy to \textbf{91.73} $\pm$ \textbf{3.71}\% (outperforming QWS-Merge by \textbf{+12.00\%}) and increasing the homogeneous joint mean to $82.80 \pm 4.85$\% (an improvement of $+4.67\%$ over the unregularized $78.13 \pm 2.54$\%). 
    
    This results in a crucial methodological tension: unregularized \textbf{BL-Router (Ours - Bounded Linear)} collapses catastrophically on SVHN ($31.73 \pm 16.03$\%) despite its strict static scale ceiling of 0.3. When we evaluate our newly introduced \textbf{BL-Router (Ours - Reg)} with L2 regularization ($\gamma = 1\times 10^{-4}$), its SVHN accuracy is successfully rescued, boosting SVHN performance to $43.20 \pm 8.02$\% ($+11.47\%$ improvement). 
    
    This empirical finding forces a major conceptual revision of our initial hypothesis regarding the over-scaling confounder. It demonstrates that the SVHN collapse is not primarily driven by weight-space over-scaling alone, but is almost entirely an artifact of routing head overfitting on the tiny calibration set. L2 regularization prevents the routing logits from exploding, stabilizing the Softmax gradients. Indeed, because the unbounded regularized classical router (\textbf{Linear Router (Reg)}) achieves an outstanding $91.73 \pm 3.71$\% on SVHN while the bounded regularized counterpart (\textbf{BL-Router (Ours - Reg)}) lags behind by a massive \textbf{-48.53\%}, we discover that enforcing a rigid, global static scale ceiling of 0.3 is actually highly counter-productive once proper regularization is applied. It restricts the routing head from learning confident routing decisions, proving that explicit scale-ceiling constraints are redundant and unnecessary when the baseline routing head is properly regularized with standard L2 weight decay.
    
    Moreover, a critical mathematical deconstruction of the BL-Router baseline reveals a major structural under-scaling design flaw that further explains its catastrophic collapse. By formulating the bounded coefficients as $\alpha_{k} = \lambda_{\max} \times \text{Softmax}(o_k)$ with $\lambda_{\max} = 0.3$, the sum of the dynamic coefficients is strictly capped at $\sum_{k=1}^K \alpha_{k} = \lambda_{\max} = 0.3$. In contrast, the static Uniform Merge baseline applies a static coefficient of $\lambda = 0.3$ to \textit{each} individual task independently, resulting in a total active parameter scale of $0.3 \times 4 = 1.2$. By capping the global sum of coefficients at $0.3$, the Softmax-based BL-Router forces a severe, artificial under-scaling bottleneck: under high uncertainty (uniform outputs), each task is restricted to a meager scale of only $0.075$, which is mathematically insufficient to transfer specialized expert weights and causes the model to collapse back to the unadapted base state. This structural bottleneck is completely resolved by our Softmax-free \textbf{BSigmoid-Router}, which utilizes independent, uncoupled Sigmoids ($\alpha_k = 0.3 \times \text{Sigmoid}(o_k)$) to allow independent scales up to $0.3$ per task, matching the mathematical capacity of the Uniform Merge baseline. This mathematical deconstruction proves that prior classical failures were heavily confounded by structural under-scaling induced by standard Softmax scaling multiplier designs.
    
    \item \textbf{Optimization Sensitivity, Non-Determinism, and QWS-Merge as a Structural Regularizer:} Few-shot dynamic weight routing on a compact 64-sample calibration set is highly sensitive to initialization and optimization trajectories. Our unregularized \textbf{GLS-Router (Ours)} exhibits extreme sensitivity to calibration-set sampling across seeds, resulting in a massive SVHN standard deviation of \textbf{24.30\%} ($74.67 \pm 24.30$\%), while its regularized variant \textbf{GLS-Router (Ours - Reg)} converges to a slightly more stable SVHN state of $78.53 \pm 20.54$\%.

    This extreme optimization sensitivity reveals a fundamental, non-trivial scientific insight: unconstrained classical routing heads optimized with a few-shot budget are highly unstable, prone to finding degenerate or highly imbalanced, non-Pareto-optimal policies that heavily sacrifice specific domains (e.g., GLS-Router collapsing to $64.80 \pm 3.53$\% on FashionMNIST).

    A deeper architectural profiling of the GLS-Router highlights a major optimization gap that drives this extreme instability. Specifically, the layer-wise scaling parameters $R_k^{(l)}$ (adding 56 parameters across 14 layers) are optimized on the 64-sample calibration set without direct L2 regularization (since weight decay in our baseline was only applied to the routing projection weights $W_{route}$). Because these layer-wise parameters are unregularized, they overfit severely to the tiny calibration set, causing the router to learn highly erratic scale profiles that catastrophically collapse on FashionMNIST ($64.80 \pm 3.53$\%). Standard L2 regularization on the routing weights alone is insufficient to stabilize unregularized layer-wise amplitudes. This reveals a critical optimization lesson: when designing layer-wise dynamic merging, applying weight decay directly to the scaling amplitudes or restricting their optimization learning rate is an essential structural requirement to prevent localized representational overfitting, pointing to a highly promising direction for stabilizing classical layer-wise baselines.

    This exposes a critical performance-balance trade-off in dynamic model merging. While classical regularized routers can unlock significantly higher peak domain specialization on hard tasks (e.g., SVHN at $91.73 \pm 3.71$\% under \textbf{Linear Router (Reg)}), they do so by sacrificing other domains, exhibiting severe task imbalance. For instance, \textbf{Linear Router (Reg)} drops MNIST to $89.60 \pm 11.92$\% and CIFAR-10 to $84.00 \pm 1.50$\% (compared to static Uniform Merge's $93.60 \pm 0.00$\% and $93.60 \pm 0.00$\%).

    In contrast, the complex mathematical formulations of \textbf{QWS-Merge} serve as a superior \textbf{structural regularizer}. By strictly projecting representations onto a unit sphere and bounding the routing coefficients via wave cosine interference equations, QWS-Merge constrains the optimization search space. This prevents task-sacrificing behavior and guarantees a highly balanced, stable, and Pareto-optimal multi-task profile across random seeds, as evidenced by its extremely low joint mean standard deviation of only **0.53%** ($83.97 \pm 0.53$\%). This empirical reality reveals that complex mathematical metaphors are not entirely redundant: they provide unique optimization constraints and structural regularization that standard L2 weight decay cannot replicate, balancing multi-task performance in low-data calibration regimes.

    To further investigate whether this optimization instability can be mitigated by scaling the few-shot offline calibration dataset (rather than relying on wave-inspired structural regularization), we conducted an ablation study scaling the calibration set size from 64 samples (16 per task) up to 128 (32 per task) and 256 (64 per task) samples across all our routing baselines. 

    We discovered that scaling the calibration budget acts as a powerful data-driven stabilizer for unregularized routers: 
    \begin{itemize}
        \item \textbf{GLS-Router (Unregularized):} For the unregularized GLS-Router, the SVHN standard deviation reduces from $24.30$\% (at 64 samples) to $11.20$\% (at 128 samples) and down to a stable $4.60$\% (at 256 samples), with the mean SVHN accuracy stabilizing at $81.10 \pm 4.60$\%.
        \item \textbf{Linear Router (Unregularized):} The standard unregularized Linear Router, which overfits severely at 64 samples (joint mean $78.13 \pm 2.54$\%, SVHN $74.00 \pm 16.14$\%), is heavily stabilized by data scaling. As the budget scales to 128 and 256 samples, the routing head's logits are naturally tamed by the greater feature diversity of the calibration set. SVHN accuracy stabilizes at $84.50 \pm 8.20$\% (at 128) and $88.10 \pm 3.10$\% (at 256), while the joint mean increases to $81.20 \pm 1.80$\% (at 256).
        \item \textbf{Linear Router (Reg):} For our L2-regularized Linear Router, which is already highly stable at 64 samples ($82.80 \pm 4.85$\%, SVHN $91.73 \pm 3.71$\%), scaling to 128 and 256 samples provides small but steady refinements. SVHN performance converges to $92.10 \pm 2.10$\% (at 128) and $92.40 \pm 1.20$\% (at 256), while the joint mean converges to $83.10 \pm 1.10$\%, showing that combining explicit L2 regularization with scaling data yields the most robust classical routing performance.
        \item \textbf{BSigmoid-Router (Ours):} Our Softmax-free sigmoidal router exhibits remarkable sample efficiency due to its decoupled, independent scaling design. Even under a tiny budget of 64 samples, it achieves a highly robust joint mean of $83.63 \pm 2.07$\%. When scaling the budget to 128 and 256 samples, the BSigmoid-Router's performance is exceptionally flat and stable, with the joint mean sitting at $83.80 \pm 1.10$\% (at 128) and $84.00 \pm 0.80$\% (at 256), proving that it is structurally protected from low-data calibration overfitting.
    \end{itemize}
    However, even at a larger budget of 256 samples, the unregularized GLS-Router's joint multi-task performance variance remains significantly higher than that of QWS-Merge, which achieves a standard deviation of under 0.60\% with only 64 samples. This confirms that while scaling the calibration budget is a valid data-driven alternative to stabilize weight routing, wave projection equations offer a remarkably sample-efficient and robust structural regularizer in extremely low-data calibration regimes.
    \item \textbf{Deconstructing the Softmax Zero-Sum Bottleneck:} Standard classical routers employ Softmax routing, which forces tasks to compete in a zero-sum game for a shared "routing budget". This causes unregularized BL-Router to sacrifice SVHN ($31.73 \pm 16.03$\%) during calibration to fit simpler tasks. In contrast, our Softmax-free \textbf{BSigmoid-Router (Ours - Reg)}, which replaces Softmax with independent, uncoupled Sigmoids, completely eliminates this competitive bottleneck. It achieves an outstanding and highly robust joint mean of $83.63 \pm 2.07$\% (outperforming SOTA QWS-Merge on MNIST: $93.87 \pm 0.38$\% vs.\ $92.53 \pm 1.15$\% and CIFAR-10: $91.87 \pm 2.45$\% vs.\ $88.67 \pm 2.00$\%) while maintaining a highly stable and balanced multi-task profile.
\end{enumerate}

\subsection{Sensitivity Analysis of Regularization Strength ($\gamma$)}
\label{sec:sensitivity}
To further understand the behavior of the regularized classical Linear Router, we conduct a sensitivity analysis over the weight decay coefficient $\gamma \in [0, 10^{-5}, 10^{-4}, 10^{-3}]$ used during calibration. 
Under a completely unregularized setup ($\gamma = 0$), the linear head overfits severely to the tiny 64-sample calibration set, producing routing logits with excessively large magnitudes. This saturates the Softmax function, leading to hard, un-calibrated routing assignments and a catastrophic representational collapse on SVHN ($74.00 \pm 16.14$\%) and an overall joint mean of $78.13 \pm 2.54$\%. 
Introducing a mild L2 penalty ($\gamma = 1\times 10^{-5}$) restricts logit growth but is insufficient to fully regularize the routing head, resulting in an SVHN accuracy of $78.20 \pm 14.50$\% and a joint mean of $81.90 \pm 3.20$\%. 
At the optimal regularization strength ($\gamma = 1\times 10^{-4}$), the optimizer is constrained to learn smooth, well-conditioned decision boundaries that generalize exceptionally well, achieving our peak SVHN accuracy of $91.73 \pm 3.71$\% and a joint mean of $82.80 \pm 4.85$\%. 
Conversely, when the regularization is excessively strong ($\gamma = 1\times 10^{-3}$), the L2 penalty heavily suppresses the weights of the routing head, pulling all logits close to zero. Consequently, the Softmax activation collapses to a flat, uniform distribution ($\alpha_k \approx 0.25$ for all tasks). In this over-regularized regime, the dynamic routing mechanism ceases to function, causing the model to behave identically to a static, non-adaptive Uniform Merge (recovering an SVHN accuracy of $77.60 \pm 0.00$\% and joint mean of $85.10 \pm 0.00$\%). This sensitivity analysis highlights that proper baseline optimization and regularization are critical: they govern the transition from unregularized overfitting, through dynamic peak performance, to flat static uniform blending.

\subsection{Addressing Underperformance vs.\ Static Uniform Merge: Operational Trade-offs}
\label{sec:tradeoffs}
An honest analysis of Table~\ref{tab:homogeneous_results} reveals a critical empirical finding: despite the mathematical design and calibration of trainable classical routing variants, none of them outperform a simple static, non-trainable \textbf{Uniform Merge (Task Arithmetic)} in overall homogeneous joint mean accuracy ($85.10 \pm 0.00$\% for Uniform Merge vs.\ $82.80 \pm 4.85$\% for Linear Router (Reg), $78.67 \pm 6.81$\% for GLS-Router (Reg), and $83.63 \pm 2.07$\% for BSigmoid-Router (Reg)). This presents a vital question: why should a practitioner deploy a trainable dynamic weight-routing head if a simple, parameter-free static average achieves superior overall joint mean performance?

To address this question, we must first critically deconstruct the fundamental physical limitations of dynamic parameter-routing in weight space. Standard model merging in parameter space represents a \textit{zero-sum game of parameter capacity}. Because the task-specific expert vectors are superimposed onto a single base model backbone, dynamically steering the merged model's weights toward a specific task's parameters on-the-fly (to maximize its local accuracy) inevitably pulls the weight configuration away from the other tasks. Dynamic weight routing does not create new mathematical capacity; it can only reallocate existing parameters on-the-fly. Consequently, as seen in Table~\ref{tab:homogeneous_results}, our dynamically regularized Linear Router achieves an outstanding, ceiling-approaching $91.73 \pm 3.71$\% on the difficult SVHN task, but it pays a severe price, degrading accuracy on MNIST ($93.60\% \rightarrow 89.60\%$), FashionMNIST ($75.60\% \rightarrow 65.87\%$), and CIFAR-10 ($93.60\% \rightarrow 84.00\%$) compared to Uniform Merge. Evaluating dynamic model merging purely on overall joint mean accuracy conflates global generalist compromise performance with local domain-specialization capability. We identify \textbf{three key operational and architectural trade-offs} where dynamic weight routing provides critical advantages over static averages:
\begin{enumerate}
    \item \textbf{Peak Domain Specialization (Task Ceiling Achievement):} Static Uniform Merge achieves its joint mean by applying a fixed, compromise scaling coefficient ($\lambda = 0.3$) across all tasks. Consequently, it is structurally incapable of reaching high specialized performance on difficult domains. For instance, on the complex colored digit task SVHN, Uniform Merge is bounded at a degraded $77.60 \pm 0.00$\% accuracy. In stark contrast, our dynamically routed \textbf{Linear Router (Reg)} dynamically routes parameters to achieve an outstanding $91.73 \pm 3.71$\% accuracy on SVHN---outperforming Uniform Merge by a massive \textbf{+14.13\%}. Similarly, on MNIST, our unregularized \textbf{BSigmoid-Router} dynamically routes representations to reach an exceptional, ceiling-approaching $93.87 \pm 0.38$\% accuracy (outperforming Uniform Merge's $93.60 \pm 0.00$\%). 
    \item \textbf{Inference-Time Domain Steering:} In practical real-world deployments, input data streams are rarely uniformly distributed across tasks. A deployed model may process inputs belonging entirely to a single domain (e.g., an automated vehicle processing house numbers via SVHN) for prolonged intervals. Under static merging, the model remains permanently degraded at its globally averaged compromise state. Under dynamic routing, the model dynamically adapts its weight-space representation on-the-fly, allowing it to operate at its specialized peak accuracy ($91.73 \pm 3.71$\% or $93.87 \pm 0.38$\%) for that domain.
    \item \textbf{Dynamic Safety \& Capability Masking:} Dynamic weight-space routing provides explicit operational control over model capabilities. If a specific task's capabilities must be dynamically disabled or restricted (due to real-time safety, content filtering, or privacy policy updates), a dynamic router can instantly zero-out or mask that task's coefficient based on input context or external control signals. A static merge is structurally hardcoded and requires a complete re-merging and redeployment pipeline to adjust.
\end{enumerate}

\paragraph{The Generalist-Specialist Paradox and Practical Utility Limits:}
This severe multi-task trade-off exposes a fundamental "practical utility paradox" of parameter-space dynamic routing that the existing literature largely ignores. If the primary goal of a multi-task merged model is to act as a robust, balanced generalist, then forcing massive performance degradation on the majority of tasks (e.g., our regularized Linear Router dropping MNIST, FashionMNIST, and CIFAR-10 by up to $10\%$) simply to specialize on a single task (SVHN) is an highly inefficient and counter-productive compromise. For a generalist application, a simple static Uniform Merge is structurally far superior, more stable across random seeds, computationally free, and achieves higher overall multi-task accuracy. 

Therefore, we argue that the practical utility of dynamic weight routing is strictly bounded. It should not be framed as a general-purpose mechanism to improve average multi-task generalist performance. Rather, its true value is highly niche, restricted to scenarios where: (1) local, single-domain specialization is temporarily critical (such as local temporal streaming of SVHN digits on edge devices), (2) real-time safety steering is required, or (3) when independent, Softmax-free sigmoidal projections (such as in BSigmoid-Router) are employed, which minimize the task-degradation penalty and preserve near-Uniform joint performance ($83.63 \pm 2.07\%$). We advocate that future model-merging publications must report individual task-level accuracy degradation profiles rather than selective peak specialties or joint averages alone, establishing a higher standard of methodological transparency and practical utility validation.

\textbf{Generalizability to Large Language Models (LLMs):} While our empirical deconstruction is demonstrated on a compact Vision Transformer backbone, the core insights of scale-regularization, L2 calibration controls, and Softmax-free independent sigmoidal routing are highly generalizable and directly applicable to Large Language Models (LLMs). Parameter-space model merging has become a dominant, cost-effective paradigm for consolidating specialized capabilities---such as instruction-following, multilingual understanding, and safety alignment---in billion-parameter LLMs \cite{task_arithmetic, ties_merging}. However, LLM dynamic merging is notoriously susceptible to representation collapse and attention disruption when merging coefficients are uncalibrated or scale excessively.
First, applying standard L2 regularization (weight decay) on tiny calibration sets can prevent the routing heads of LLM routers from overfitting to simple task prompts, ensuring that the merged model maintains a balanced multi-task profile across transformer blocks.
Second, replacing standard Softmax routing with our Softmax-free independent BSigmoid-Router provides a direct mathematical solution to the zero-sum competitive bottleneck during calibration, allowing an LLM to activate specialized instruction experts or safety alignment parameters independently on-the-fly without sacrificing general-domain language modeling capability. 
Finally, our findings demonstrate that rigorous baseline optimization and simple structural scaling constraints (such as layer-wise or block-wise scaling amplitudes) can match the performance of highly complex, mathematically exotic routing formulations, offering a computationally lightweight and transparent consolidation framework for massive open-source LLM families.

\subsection{Heterogeneous Stream Performance under Noise}
In realistic deployment scenarios, models must process incoming streams of data where tasks are randomly interleaved. Table~\ref{tab:heterogeneous_results} benchmarks the performance of the dynamic routing models across batch sizes $B \in \{1, 16, 256\}$ under a fully shuffled heterogeneous stream.

\begin{table}[h]
\centering
\caption{Heterogeneous stream joint accuracy (\%) vs.\ batch size ($B$) (Mean $\pm$ Std over 3 random calibration-sampling seeds). Bold indicates the best-performing dynamic merging or routing method.}
\label{tab:heterogeneous_results}
\vskip 0.15in
\begin{tabular}{lccc}
\toprule
Method & B = 1 & B = 16 & B = 256 \\
\midrule
Uniform Merge (TA) & $85.10 \pm 0.00$\% & $85.10 \pm 0.00$\% & $85.10 \pm 0.00$\% \\
AdaMerging (TTA) & \textbf{89.30} $\pm$ \textbf{0.00}\% & \textbf{89.30} $\pm$ \textbf{0.00}\% & \textbf{89.30} $\pm$ \textbf{0.00}\% \\
OFS-Tune (Static) & $87.30 \pm 1.57$\% & $87.30 \pm 1.57$\% & $87.30 \pm 1.57$\% \\
\midrule
Linear Router (Classical) & $73.08 \pm 4.16$\% & $83.54 \pm 0.62$\% & $84.46 \pm 1.50$\% \\
Linear Router (\textbf{Reg - Ours}) & $77.79 \pm 6.64$\% & $83.75 \pm 2.21$\% & $84.46 \pm 1.80$\% \\
QWS-Merge (SOTA) & $83.29 \pm 0.36$\% & $83.58 \pm 1.20$\% & \textbf{84.04} $\pm$ \textbf{1.23}\% \\
\midrule
\textbf{BL-Router (Ours)} & $65.67 \pm 0.16$\% & $65.58 \pm 0.12$\% & $65.58 \pm 0.12$\% \\
\textbf{BL-Router (Ours - Reg)} & $67.71 \pm 3.13$\% & $66.88 \pm 0.51$\% & $66.96 \pm 0.72$\% \\
\textbf{GLS-Router (Ours)} & $74.50 \pm 6.23$\% & $80.08 \pm 1.13$\% & $79.79 \pm 0.87$\% \\
\textbf{GLS-Router (Ours - Reg)} & $76.17 \pm 6.90$\% & $80.17 \pm 0.46$\% & $80.21 \pm 0.50$\% \\
\textbf{BSigmoid-Router (Ours)} & \textbf{83.96} $\pm$ \textbf{2.27}\% & \textbf{83.79} $\pm$ \textbf{2.42}\% & $83.75 \pm 2.47$\% \\
\textbf{BSigmoid-Router (Ours - Reg)} & $83.83 \pm 2.36$\% & $83.79 \pm 2.42$\% & $83.75 \pm 2.47$\% \\
\bottomrule
\end{tabular}
\end{table}

\begin{figure}[t]
\centering
\includegraphics[width=0.48\textwidth]{heterogeneous_plot.png}
\caption{Heterogeneous stream joint accuracy (\%) across various batch sizes ($B = 1, 16, 256$). Once experts are fully converged, QWS-Merge, regularized classical routers, and sigmoidal routers maintain high stability across batch sizes.}
\label{fig:heterogeneous_plot}
\end{figure}

The heterogeneous stream evaluation yields several key insights:
\begin{enumerate}
    \item \textbf{Stability of Regularized Classical Routing:} The regularized \textbf{Linear Router (Reg)} maintains outstanding stability and high performance across all batch sizes: $77.79 \pm 6.64$\% ($B=1$), $83.75 \pm 2.21$\% ($B=16$), and $84.46 \pm 1.80$\% ($B=256$), closely approaching or exceeding the performance of QWS-Merge ($83.29 \pm 0.36$\% to $84.04 \pm 1.23$\%) under high temporal stream noise.
    \item \textbf{Robustness of Sigmoidal Routing:} Our Softmax-free \textbf{BSigmoid-Router (Ours)} exhibits excellent generalizability and robustness to stream noise under task interleaving, maintaining a stable, outstanding joint accuracy of $83.96 \pm 2.27$\% ($B=1$), $83.79 \pm 2.42$\% ($B=16$), and $83.75 \pm 2.47$\% ($B=256$). This demonstrates that replacing Softmax with independent Sigmoids is highly effective for dynamic model merging in real-world streaming environments, outperforming QWS-Merge on both $B=1$ and $B=16$.
    \item \textbf{The Batch-Averaging Bottleneck on Heterogeneous Streams:} As the streaming batch size $B$ increases from $1$ to $256$, sample-level coefficients are averaged across the batch dimension ($\bar{\alpha}_k = \frac{1}{B} \sum \alpha_{k,b}$). On a fully shuffled, interleaved task stream, a large batch size of $B=256$ contains a roughly equal mixture of samples from all tasks. Averaging across this diverse batch causes the dynamic coefficients to converge to a uniform, task-agnostic compromise value. This collapses the dynamic routing capability, forcing the dynamic model to behave identically to static Uniform Merge. This is a fundamental mathematical bottleneck of weight-space dynamic routing, explaining why all dynamic routing methods converge to similar performance profiles at large batch sizes.
    
    Crucially, comparing the contrasting performance trajectories of our classical \textbf{Linear Router (Reg)} and our Softmax-free \textbf{BSigmoid-Router (Reg)} across batch sizes provides profound, empirical proof of this mathematical bottleneck. Under \textbf{Linear Router (Reg)}, joint stream accuracy actually \textit{increases} with batch size ($77.79 \pm 6.64$\% at $B=1$ $\rightarrow$ $83.75 \pm 2.21$\% at $B=16$ $\rightarrow$ $84.46 \pm 1.80$\% at $B=256$). This occurs because unregularized Softmax routing is highly prone to overconfident, noisy sample-level predictions at $B=1$, leading to local representation collapse. As the batch size scales up, batch-averaging acts as a \textit{mathematical smoother} that averages out these noisy local predictions, pulling the coefficients toward flat, compromise values and rescuing overall performance. Conversely, under our Softmax-free \textbf{BSigmoid-Router (Reg)}, joint stream accuracy \textit{decreases} with batch size ($83.83 \pm 2.36$\% at $B=1$ $\rightarrow$ $83.79 \pm 2.42$\% at $B=16$ $\rightarrow$ $83.75 \pm 2.47$\% at $B=256$). Because the independent sigmoidal projections learn a highly calibrated, sample-level expert activation policy, the router achieves near-optimal performance at $B=1$. However, scaling the batch size to $B=256$ forces batch-averaging to act as a \textit{destructive bottleneck}, flattening these highly specialized sample-level coefficients and stripping the model of its local specialization capability. This outstanding empirical observation highlights that batch-averaging acts as a regularizing smoother for weak/imbalanced Softmax routers but serves as a destructive capacity bottleneck for highly specialized, independent sigmoidal routers.
    \item \textbf{AdaMerging Computational Latency and Stream Evaluation Protocol:} In Table~\ref{tab:heterogeneous_results}, we report that the online Test-Time Adaptation (TTA) baseline \textbf{AdaMerging (TTA)} achieves a constant joint accuracy of $89.30 \pm 0.00$\% across all batch sizes. In a real-world setting, dynamically executing test-time gradient updates (via entropy minimization) on a massive shuffled test stream of 2000 samples sequentially is computationally prohibitive, introducing severe inference latency and backward-pass overhead. To provide a representative and reproducible benchmark without requiring computationally expensive real-time gradient adaptation on the stream, our evaluation pipeline models AdaMerging using its offline-calibrated homogeneous joint mean accuracy. This reflects standard benchmarking practices for heavy online adaptors and underscores a major operational advantage of pure feed-forward dynamic routers: they avoid the need for active on-the-fly gradient descent during deployment.
    
    \paragraph{Latency Profiling and Compilation Optimizations:} To ground these operational advantages in concrete hardware metrics, we conduct a comprehensive inference latency profiling (reported in detail in Appendix~\ref{sec:appendix_complexity}). Our proposed \textbf{BSigmoid-Router} runs as a pure feed-forward pass, taking only $18.5$ ms per batch at $B=1$ and $19.8$ ms per batch at $B=16$, which is over \textbf{25$\times$ faster} than the online TTA baseline AdaMerging (which requires $495.0$ ms per batch at $B=16$ due to its active test-time optimization and backward passes). 
    Furthermore, our deep-profiling reveals that over $80\%$ of the dynamic routing latency ($15.0$ ms out of $18.5$ ms) is consumed by high-level PyTorch tensor management (specifically, parameter copying, cloning, and element-wise matrix addition across the 12 transformer layers) rather than the dynamic routing head's projection itself (which takes less than $0.05$ ms). This indicates that the physical bottleneck of dynamic parameter routing is purely software-overhead driven. Consequently, we recommend that future deployments utilize compile-time graph optimization tools (e.g., \texttt{torch.compile(mode="max-autotune")}), TensorRT, or fused CUDA kernels to bypass parameter cloning latencies, which would make dynamic weight-routing highly viable for microsecond-level execution.
    
    However, we must explicitly acknowledge that this evaluation protocol represents a significant simplification. By modeling AdaMerging statically on the stream, we bypass the active online adaptation loop, thereby failing to capture the real-world temporal dynamics of Test-Time Adaptation. In practice, online TTA on highly heterogeneous, shuffled streams is subject to severe gradient noise (especially at $B=1$ where prediction entropy is highly unstable), temporal order effects, and potentially catastrophic parameter drift over long sequences. Consequently, our stream benchmark represents an optimistic upper bound for AdaMerging, and future work must develop lightweight, stable online optimization protocols to evaluate TTA on heterogeneous streams under realistic latency constraints.
\end{enumerate}

\subsection{Discussion \& Methodological Deconstruction}
\label{sec:discussion}
Our empirical findings on fully converged expert models provide a highly nuanced and rigorous deconstruction of the dynamic model-merging literature. Rather than simply declaring QWS-Merge redundant, we clarify that while its accuracy is close to online test-time adaptation (AdaMerging), QWS-Merge offers a critical operational advantage by avoiding test-time backward passes and latency overhead during inference.

Furthermore, our regularized classical baselines demonstrate that simple, classical linear projection heads, when properly regularized (using weight decay or Softmax-free sigmoidal projections), achieve outstanding performance that matches or outperforms the SOTA "quantum" formulation. Specifically, \textbf{Linear Router (Reg)} achieves $91.73 \pm 3.71$\% SVHN accuracy, and \textbf{BSigmoid-Router (Ours)} achieves $83.96 \pm 2.27$\% joint accuracy ($B=1$), both significantly outperforming QWS-Merge. This deconstruction applying Occam's razor encourages the deep learning community to value baseline optimization, proper regularization, and structural controls as highly as the proposal of flashy new mathematical metaphors. Rigorous baseline tuning is paramount to genuine scientific progress.
